\section{Global preparation and finite assembly}
\label{sec:global-construction}

We now arrange a finite sequence of primitive additions to which
Proposition~\ref{prop:primitive-addition} applies.  The main point is the order
of the choices.  The initial immersion will be modified at the crossings of
the support boundaries, although those boundaries depend on the number of
primitive additions.  Uniform second-derivative bounds for the initial
modification allow us to choose the number of additions first.

The assertion is immediate if $M$ is empty.  The connected components
of a smooth manifold are open.  Compactness
therefore implies that $M$ has only finitely many components.  It suffices
to carry out the construction on each component: the resulting maps
together define one smooth isometric immersion on $M$.  We may thus
assume in this section that $M$ is connected.

\subsection{A spherical starting immersion}

Every closed smooth surface admits an ordinary smooth immersion into $S^3$.
Indeed, the two-dimensional case of Whitney's immersion theorem gives a
$C^1$ immersion into $\R^3$ \cite[Theorem~8, p.~274]{Whitney1944}.  Approximation in $C^1$ by a
smooth map preserves injectivity of the differential on a compact domain.
Such an approximation can be obtained by mollifying in finitely many charts
and combining the approximations with a smooth partition of unity.  Finally,
inverse stereographic projection is a diffeomorphism from $\R^3$ onto a
punctured $S^3$.  Neither step requires an orientation of the surface.

For an immersion $I:M\to S^3$, denote its second fundamental form in $S^3$
by $B_I^{S^3}$, regarded as an $\R^4$-valued tensor.  This notation avoids any
choice of unit normal to the surface in $S^3$.

\begin{lemma}[Finite-point spherical preparation]
\label{lem:spherical-preparation}
Let $M$ be a closed smooth surface and let $I:M\to S^3$ be a smooth
immersion.  There are constants $C,c>0$ and $\varepsilon_0>0$ with the
following property.  For every finite set $P\subset M$ and every
$0<\varepsilon<\varepsilon_0$, there is a smooth immersion
$I_P:M\to S^3$ such that
\[
 B_{I_P}^{S^3}(p)=0\quad(p\in P),\qquad
 \|I_P-I\|_{C^1}<\varepsilon,\qquad
 \|I_P\|_{C^2}\le C,
 \qquad I_P^*\delta_4\ge c\,g_{\mathrm{bg}},
\]
where $g_{\mathrm{bg}}$ is the fixed background metric.  The constants are
independent of the number and separation of the points of $P$.
\end{lemma}

\begin{proof}
Choose a finite atlas with precompact smaller charts covering $M$.  For each
$p\in P$, use one of these charts, translated so that $p$ has coordinates
$u=0$.  Choose radii $r_p>0$ so that the corresponding coordinate balls have
pairwise disjoint closures and lie in the larger charts.  Let $\chi$ be a
fixed smooth bump, equal to one near zero and supported in the unit ball.
In the chart at $p$, put
\[
 Q_p(u)=-\frac12\chi(u/r_p)
       \sum_{i,j=1}^2 B_I^{S^3}(\partial_i,\partial_j)(p)u_i u_j,
\]
and extend this function by zero.  Set
\[
 J=I+\sum_{p\in P}Q_p,
 \qquad I_P=\frac{J}{|J|}.
\]
For sufficiently small radii, $|J|$ stays uniformly away from zero.

At $p$, the value and first derivative of $Q_p$ vanish and
$\partial_i\partial_jQ_p=-B_I^{S^3}(\partial_i,\partial_j)(p)$.
The latter vector is tangent to $S^3$ at $I(p)$ and normal to $dI(T_pM)$.
The derivative of $v\mapsto v/|v|$ at $I(p)$ is projection to
$T_{I(p)}S^3$.  Comparing its second-derivative chain rule for $J$ and for
$I$ therefore gives
\[
 dI_P(p)=dI(p),\qquad
 \partial_i\partial_j I_P(p)
 =\partial_i\partial_j I(p)
   -B_I^{S^3}(\partial_i,\partial_j)(p).
\]
Thus the spherical normal part of the second derivative vanishes at $p$.

The $C^0,C^1,C^2$ norms of $Q_p$ are bounded respectively by
$C_0r_p^2,C_0r_p,C_0$, with $C_0$ uniform in $p$.  Since the supports are
disjoint, the corresponding norms of the sum are bounded by the maximum of
these bounds, up to the fixed atlas constants.  Normalization has uniformly
bounded derivatives on the range in question.  Taking the radii small gives
the prescribed $C^1$ closeness and a uniform $C^2$ bound.  Compactness and
the immersion property of $I$ then give the metric lower bound and preserve
immersion, after fixing $\varepsilon_0$ sufficiently small.
\end{proof}

Only the first two derivatives are controlled uniformly in this lemma.
Higher derivatives may grow when the points are close together.  The
preparation below uses only the stated bounds; constants requiring higher
derivatives will be chosen after the finite set has been fixed.

For later use, consider $F^0=\lambda I_P$, where $\lambda>0$, and write
$\gamma^0=\gamma(F^0)$.  The vector field $n^0=-I_P$ is a global unit normal
in $\R^4$.  The second fundamental forms satisfy
\begin{equation}
 B_{F^0}(v,w)
   =\lambda B_{I_P}^{S^3}(v,w)
       +\lambda^{-1}\gamma^0(v,w)n^0.
 \label{eq:spherical-radial-form}
\end{equation}
In particular, $\langle B_{F^0}(v,v),n^0\rangle>0$ for $v\ne0$.
At each point of $P$, the first term vanishes and the Gauss equation gives
$K_{\gamma^0}=\lambda^{-2}$.  For any nonzero vectors $Y_i,Y_l$ at such a
point, the quantity in~\eqref{eq:boundary-invariants} is consequently
\begin{equation}
 Q_{il}
 =\lambda^{-2}\bigl(
      \gamma^0(Y_i,Y_l)^2+|Y_i\wedge Y_l|_{\gamma^0}^2\bigr)
 =\lambda^{-2}\gamma^0(Y_i,Y_i)\gamma^0(Y_l,Y_l)>0.
 \label{eq:initial-crossing-positivity}
\end{equation}
This identity also holds when the two vectors are proportional.

\subsection{A finite list with prescribed intermediate metrics}

\begin{lemma}[Global primitive preparation]
\label{lem:global-preparation}
Let $(M,g)$ be a closed smooth Riemannian surface.  There exist a smooth
immersion $F^0:M\to\R^4$, a smooth unit normal $n^0$, and a finite list
\[
 (D_j,x_j,y_j,a_j),\qquad 1\le j\le K,
\]
with the following properties.  Each $D_j$ is a coordinate disk with smooth
boundary $C_j$; $(x_j,y_j)$ are coordinates on a neighborhood of
$\overline D_j$; and $a_j$ is smooth, positive on $D_j$, and zero on its
exterior.  Set $Y_j=\partial_{y_j}$ and
\[
 \gamma^j=\gamma(F^0)+\sum_{i=1}^j a_i^2(dx_i)^2,
 \qquad \gamma^0=\gamma(F^0).
\]
Then:
\begin{enumerate}
\item $\gamma^K=g$, and
      $\Hess_{\gamma^{j-1}}x_j(Y_j,Y_j)>0$ on $\overline D_j$.
\item The boundaries are pairwise transverse and have no triple
      intersection.  For $i\ne j$, $Y_j$ is tangent to $C_i$ at only finitely
      many points of $\overline D_j$.  At $C_i\cap C_l\cap D_j$, with
      $i,l,j$ distinct, both $dx_j(Y_i)$ and $dx_j(Y_l)$ are nonzero.
\item $B_{F^0}$ is nowhere the zero tensor, and $F^0,n^0$ satisfy all
      boundary conditions~\eqref{eq:boundary-invariants} for this list.
\end{enumerate}
\end{lemma}

\begin{proof}
Fix a spherical immersion $I$ and the uniform family of modifications in
Lemma~\ref{lem:spherical-preparation}, without yet specifying the finite set.
We make the choices below uniformly for this family.

\paragraph{Metric and cone margins.}
At any point, take a $g$-orthonormal coframe $e^1,e^2$ and the covectors
\[
 \ell_1=e^1,\qquad
 \ell_2=-\tfrac12e^1+\tfrac{\sqrt3}{2}e^2,\qquad
 \ell_3=-\tfrac12e^1-\tfrac{\sqrt3}{2}e^2.
\]
Their squares form a basis of symmetric tensors and
$g=\frac23\sum_r\ell_r^2$ at that point.  On the compact sets in the fixed
atlas, these triples have uniform bounds, inverse-basis bounds, and a
uniform neighborhood in which all three coefficients remain positive.
Choose $\lambda>0$ sufficiently small that, for every allowed modification
$\bar I$, the metrics
\[
 \gamma^0_{\bar I}=\lambda^2\bar I^*\delta_4,
 \qquad h_{\bar I}=g-\gamma^0_{\bar I}
\]
are positive, and $h_{\bar I}$ lies within this fixed cone margin around
$g$.  This one scaling works for all centers subsequently chosen.
The paths
\[
 \gamma_{\bar I,s}=\gamma^0_{\bar I}+s h_{\bar I},
 \qquad 0\le s\le1,
\]
have uniform positive metric lower bounds and uniform $C^1$ bounds.  Their
Christoffel symbols in the fixed atlas are therefore uniformly bounded.

\paragraph{Convex phases and disks.}
In background coordinates $u$ centered at a disk center, replace the linear
phases by
\[
 x_r(u)=\ell_r\cdot u+\tfrac12L|u|^2.
\]
For every metric in the above family,
\[
 (\Hess_\gamma x_r)_{ab}
   =L\delta_{ab}-\Gamma^c_{ab}(\gamma)(\ell_{r,c}+Lu_c).
\]
Choose $L$ large using the uniform connection and covector bounds, and then
choose radii so small that $L|u|$ is small on the enlarged disks.  The
displayed Hessians are uniformly positive definite there.  Shrinking
further preserves the positive basis decompositions of $h_{\bar I}$ and
the nonvanishing of $dx_r$.  A linear complementary coordinate $y_r$ makes
$(x_r,y_r)$ a chart: along each line slice with fixed $y_r$, the derivative
of $x_r$ in the $\ell_r$ direction is strictly positive.  Take finitely many
such disks $D_\nu$ covering $M$, with their closures inside the enlarged
charts.

Let $q_{\nu r}^0$ be the smooth linear coefficient forms determined by the
three unperturbed phase squares on the enlarged disk.  Thus
\[
 H=\sum_{r=1}^3q_{\nu r}^0(H)(dx_{\nu r})^2
\]
for every symmetric tensor there, and the coefficients of every
$h_{\bar I}$ have a uniform positive lower bound.  Choose smooth nonnegative
functions $\psi_\nu$, positive exactly on $D_\nu$, with
$\sum_\nu\psi_\nu^2=1$.  To obtain them, take flat positive bumps on the
disks and divide by the square root of the sum of their squares.  The sum
is positive because the disks cover $M$.

\paragraph{Perturbed cycles and their exact amplitudes.}
One cycle consists of the finite disk/triple list.  Permit small independent
changes of each disk radius and of each phase's linear part, leaving its
quadratic part fixed.  A small dilation gives a cutoff $\psi_j$ positive
exactly on the changed disk and close to the old cutoff in $C^1$.  Restrict
the perturbations so that the charts, cone margins and positive phase
Hessians for all path metrics retain uniform strict margins.  For each
perturbed cycle, indexed by $c$ starting at zero, define its own bundle endomorphism of
$\operatorname{Sym}^2T^*M$ by
\[
 \mathcal A_c(H)=\sum_{j=(\nu,r)\text{ in cycle }c}
       \psi_j^2 q_{\nu r}^0(H)(dx_j)^2.
\]
The terms extend smoothly by zero.  At zero perturbation, $\mathcal A_c=\Id$.
Hence, in a sufficiently small fixed perturbation neighborhood,
$\mathcal A_c$ is invertible with a uniformly bounded $C^1$ inverse.
For example, the derivative bound follows in local frames from
$d(\mathcal A_c^{-1})=-\mathcal A_c^{-1}(d\mathcal A_c)\mathcal A_c^{-1}$.
Shrinking that neighborhood preserves a uniform positive lower bound for
$q_{\nu r}^0(\mathcal A_c^{-1}h_{\bar I})$.

If there are $N_c$ cycles, define the amplitudes for $j$ in cycle $c$ by
\begin{equation}
 a_j=N_c^{-1/2}\psi_j
          \sqrt{q_{\nu r}^0(\mathcal A_c^{-1}h_{\bar I})}.
 \label{eq:global-amplitudes}
\end{equation}
The square root is smooth on a neighborhood of the cutoff support.  Thus
$a_j$ extends smoothly by zero and is positive exactly on its disk.  The
defining endomorphism gives the exact identity
\begin{equation}
 \sum_{j\text{ in cycle }c}a_j^2(dx_j)^2
 =N_c^{-1}\mathcal A_c(\mathcal A_c^{-1}h_{\bar I})
 =N_c^{-1}h_{\bar I}.
 \label{eq:cycle-identity}
\end{equation}

Every tensor on the left has $C^1$ norm at most $C/N_c$, uniformly over
allowed starting modifications and all permitted perturbations.  The
number of terms in one cycle is fixed.  Consequently every partial cycle
differs from its starting path metric
$\gamma^0_{\bar I}+(c/N_c)h_{\bar I}$ by at most $C'/N_c$ in $C^1$, where
$c$ is the number of completed cycles.  All partial metrics are uniformly
positive.  On this bounded set, the Christoffel symbols are Lipschitz
functions of the metric coefficients and their first derivatives.
Choose $N_c$ large enough that the Hessian change caused by a partial cycle
is less than half the fixed positive margin.  Every phase then has positive
Hessian for the metric immediately before its own addition.  This choice
does not depend on the crossing set.

At this stage the number of terms and the admissible perturbation
neighborhoods are fixed.  It remains to choose the perturbations and the
particular starting map; all of the estimates just proved hold for those
choices.

\paragraph{Boundary and phase transversality.}
Write $\rho_i$ for the squared coordinate radius defining the $i$th disk.
For each pair of disks, apply Sard's theorem
\cite[Theorem~4.1, p.~885]{Sard1942} to $(\rho_i,\rho_l)$ on the overlap of
their enlarged charts.  Outside a null set of level pairs, the boundary
curves intersect transversely.  For each triple, the image of
$(\rho_i,\rho_l,\rho_j)$ has three-dimensional measure zero.
Indeed, on a compact coordinate box this smooth map has bounded
derivative; a cover by $O(\epsilon^{-2})$ squares of side $\epsilon$
gives an image cover by cubes of total volume $O(\epsilon)$.
Countable compact exhaustion proves the assertion on the overlap.
Squaring the positive radii is a smooth change of parameters with a
locally Lipschitz inverse.  Each exceptional set remains null when
lifted to the finite product of the other allowed radius intervals.
There are only finitely many pairs and triples, so their exceptions
can be avoided simultaneously in every permitted open parameter box.
Repeated copies of the same disk cause no difficulty: distinct nearby
radii give disjoint boundaries.  The resulting boundaries are compact
embedded circles, and their transverse pairwise intersections are finite.

Keep these boundaries fixed and vary the phase linear parts.  On a regular
parametrized boundary arc $u(t)$ in the enlarged phase chart $U_j$, put
\[
 H(t,\ell)=\frac{d}{dt}x_\ell(u(t))
           =(\ell+Lu(t))\cdot u'(t).
\]
Since $\partial_\ell H=u'(t)\ne0$, the set $Z=\{H=0\}$ is a smooth
two-dimensional surface in the three variables $(t,\ell)$.
Its tangent equation
\[
 H_t\,\dot t+\partial_\ell H\cdot\dot\ell=0
\]
shows that the projection $Z\to\R^2_\ell$ is singular precisely when
$H_t=0$.  Sard's theorem in countably many coordinate charts of $Z$
therefore excludes a null set of linear-phase parameters, outside which
every critical point of the restricted phase is nondegenerate.
Use a countable collection of boundary arcs covering the entire open
set $C_i\cap U_j$.  The set $C_i\cap\overline D_j$ is compact and lies
inside $U_j$; an infinite sequence of critical points there would
accumulate at another critical point and contradict its nondegeneracy.
Thus there are only finitely many tangencies of $\ker dx_j$, equivalently
$Y_j$, to $C_i$ on this compact portion.

At each already fixed crossing $p\in C_i\cap C_l\cap D_j$, require
$dx_j(p)$ to be nonproportional to both $dx_i(p)$ and $dx_l(p)$.
Holding the other phase parameters fixed, each forbidden choice of
$\ell_j$ lies on a proper affine line: varying $\ell_j$ varies
$dx_j(p)$ arbitrarily, and all phase covectors remain nonzero in the
allowed neighborhood.  Fubini's theorem makes these exceptions null
in the full phase parameter space.  The preceding single-phase
exceptions remain null there as well.  Since there are finitely many
index pairs and crossings, all these requirements can be satisfied
simultaneously in every permitted open parameter box.
In two dimensions the determinant conditions say exactly
$dx_j(Y_i)\ne0$ and $dx_j(Y_l)\ne0$ at these crossings.

\paragraph{The final choice of the starting map.}
Let $P$ be the finite set of all boundary crossings just selected.  Apply
Lemma~\ref{lem:spherical-preparation} to this set, and set
$F^0=\lambda I_P$, $n^0=-I_P$ and $h=g-\gamma(F^0)$.
Use this actual $h$ in~\eqref{eq:global-amplitudes}.  The uniform family
bounds ensure that all previous cone, cycle and Hessian estimates remain
valid.  Summing~\eqref{eq:cycle-identity} over the $N_c$ cycles gives
$\gamma^K=\gamma^0+h=g$.

Equation~\eqref{eq:spherical-radial-form} gives a positive dot product
with $n^0$ for every pure second-form direction, everywhere on $M$.
It implies nonvanishing of $B_{F^0}$ and the non-antipodality conditions
on all the boundaries.  At every crossing,
\eqref{eq:initial-crossing-positivity} gives both ordered inequalities
$Q_{il}>0$ and $Q_{li}>0$.  This proves all the assertions.
\end{proof}

\subsection{Completion of the global construction}

\begin{proof}[Proof of Theorem~\ref{thm:main}]
Take the finite list supplied by Lemma~\ref{lem:global-preparation}.
Starting from $F^0,n^0$, apply Proposition~\ref{prop:primitive-addition}
to $D_1,x_1,y_1,a_1$, with the remaining boundaries as the later curves.
The proposition produces a smooth immersion with metric $\gamma^1$, a
global preferred unit normal, a nowhere-zero second-form tensor, and the
boundary conditions for all remaining indices.

Inductively, before step $j$, the current metric is exactly
$\gamma^{j-1}$.  Its required phase-Hessian inequality was established in
Lemma~\ref{lem:global-preparation}; the curve and direction conditions are
fixed throughout the construction.  The previous primitive step supplies
the other hypotheses.  Thus the proposition applies at each of the
finitely many steps.  Its constants and oscillation parameters are chosen
for the current fixed smooth data; no uniform higher-derivative bounds
over different steps are needed.

The resulting smooth map $F^K$ has
\[
 (F^K)^*\delta_4=\gamma^K=g.
\]
Since $g$ is positive definite, $dF^K$ is injective at every point.
This is the required isometric immersion.

All maps and corrections were constructed on $M$.  The starting normal
was the radial section $-I_P$, and Proposition~\ref{prop:primitive-addition}
propagates a single global normal section using local disk constructions
and exterior $C^2$ closeness.  It does not require the new map to equal the
old map on the exterior.  Neither a global oriented normal frame nor an
orientation of $M$ was used, so the argument includes nonorientable
surfaces.
\end{proof}
